% A.107: proved comparison maps; an explicit arithmetic obstruction is supplied in A134/A137.
\subsection{The exact index-two comparison for the harmonic puncture}

\begin{statement}{A.107}{Proved}
{The harmonic cover, its exact Selmer fibres, and finite quotient maps}
Let
$V=\mathbf P^1_\Q\setminus\{0,1,\tfrac12,\infty\}$ and
$U=\mathbf P^1_\Q\setminus\{0,1,\infty\}$.
The map $t=4T(1-T)$ is a finite etale cover $f:V\to U$
of degree two, with deck involution $T\mapsto1-T$.
For every completed global triple
$(A_0,A_1,A_{1/2})$ whose localization at each finite
prime is the triple of an ordinary point of $V$, put
$A=-4A_0A_1$ and $B=4A_{1/2}^2$.
Pushforward gives an exact bijection
\[
 \operatorname{Sec}_{\rm Sel}(V;A_0,A_1,A_{1/2})
       \xrightarrow{\ \sim\ }
 \operatorname{Sec}_{\rm Sel}(U;A,B).
\]
The real Selmer condition is retained.
The two possible global lifts before fixing the old
coordinates differ by the single deck involution,
which sends the old triple to
$(-A_1,-A_0,-A_{1/2})$.

For compatible rational base points write
$\Lambda=\pi_1(V_{\overline\Q})\triangleleft
\Pi=\pi_1(U_{\overline\Q})$ and set
$H_4=\ker(\Lambda\to\Lambda^{\rm ab}/4)$,
$K_4=\overline{H_4^2[H_4,H_4]}$.
Then $K_4$ is normal in $\Pi$ and is stable under the
point-induced arithmetic action.  There are actual
arithmetic quotient maps with geometric sequence
\[
 \begin{gathered}
 1\longrightarrow\Lambda/K_4\longrightarrow\Pi/K_4
       \longrightarrow C_2\longrightarrow1,\\
 |\Lambda/K_4|=2^{135},\qquad|\Pi/K_4|=2^{136}.
 \end{gathered}
\]
Both the third-derived quotient $\Pi/\Pi'''$ and
the lower-central quotient $\Pi/\Gamma_{136}\Pi$
map onto $\Pi/K_4$, compatibly with the arithmetic
extensions.  These assertions prove the stated
comparison maps without asserting the existence of
an obstructed harmonic completed triple.
\end{statement}
\begin{proof}
Put
\[
 V=\mathbf P^1_\Q\setminus\{0,1,\tfrac12,\infty\},
 \qquad U=\mathbf P^1_\Q\setminus\{0,1,\infty\}.
\]
The rational morphism and its deck involution are
\[
 f:V\longrightarrow U,\qquad t=4T(1-T),
 \qquad \iota(T)=1-T,\qquad
                   1-t=(1-2T)^2.
 \tag{A.107.1}
\]
The derivative is $4(1-2T)$; the only ramification points
on the proper projective lines are $T=1/2$ and $\infty$,
both of index two.
Their images are $1$ and $\infty$, and the inverse image
of $0$ is $\{0,1\}$.
Thus (A.107.1) restricts to the actual finite etale cover of
degree two displayed above.
Its root coordinate is $S=1-2T$ and its inverse expression
is $T=(1-S)/2$.

\paragraph{Arithmetic groups and the Kummer morphisms.}
Choose $R\in V(\Q)$ and use the compatible geometric
base points at $R$ and $t_0=f(R)$.
Write $\Pi=\pi_1(U_{\overline\Q})$ and
$\Lambda=\pi_1(V_{\overline\Q})$.
The finite-cover correspondence, with geometric base change
as in \cite[Expos\'e IX, Th\'eor\`eme 6.1]{AgentSGA1},
gives the literal inclusions and quotient
\[
 \pi_1(V,R)\lhook\joinrel\longrightarrow\pi_1(U,t_0),
 \qquad
 \Lambda=\ker\{\Pi\longrightarrow\mu_2\},
 \qquad [\Pi:\Lambda]=2.
 \tag{A.107.2}
\]
The quotient character is that of $1-t$.
Since $1-t_0=(1-2R)^2$ is a rational square, the
point-induced splitting identifies this arithmetic
quotient with $\mu_2\times G_\Q$.
There is no nonzero coboundary for its trivial
$G_\Q$-module $\mu_2$.
Consequently a section whose Kummer class of $1-t$ is
a square has zero literal quotient character and lies
in the subgroup $\pi_1(V,R)$.

The three old coordinate functions and the two new
ones satisfy the exact identities
\[
 \begin{aligned}
 A_0&=\kappa(T),&
 A_1&=\kappa(T-1),&
 A_{1/2}&=\kappa(T-\tfrac12),\\
 A&=-4A_0A_1,&
 B&=4A_{1/2}^{\,2}.
 \end{aligned}
 \tag{A.107.3}
\]
Here constants denote their actual Kummer classes, and
the same formulas hold after pulling back along any
section.  They are induced by (A.107.1), rather than an
identification between unspecified quotient groups.
The deck involution acts on the three old classes by
\[
               (A_0,A_1,A_{1/2})
                    \longmapsto(-A_1,-A_0,-A_{1/2}).
 \tag{A.107.4}
\]
Both old triples in (A.107.4) have exactly the same new
pair under (A.107.3).

\paragraph{The exact comparison of Selmer-section fibres.}
Let $(A_0,A_1,A_{1/2})$ be any completed global triple
with ordinary finite-place coordinates: at each finite
prime $p$ there is $T_p\in V(\Q_p)$ with
\[
 (A_0)_p=\kappa(T_p),\quad
 (A_1)_p=\kappa(T_p-1),\quad
 (A_{1/2})_p=\kappa(T_p-\tfrac12).
 \tag{A.107.5}
\]
Define $(A,B)$ by (A.107.3).
It is an ordinary completed unit-equation pair, with
the actual finite-place centre $t_p=f(T_p)$.
The pushforward of sections along (A.107.2) induces a
bijection
\[
 \operatorname{Sec}_{\rm Sel}
       \bigl(V;A_0,A_1,A_{1/2}\bigr)
       \xrightarrow{\ \sim\ }
 \operatorname{Sec}_{\rm Sel}\bigl(U;A,B\bigr).
 \tag{A.107.6}
\]
These are the usual geometric-conjugacy classes, with
the indicated complete Kummer coordinates.  The real
condition in the Selmer definition is retained on both sides.

To prove surjectivity, take a section in the right
side.  At a finite prime its ordinary pair is
$(\kappa(t_p),\kappa(1-t_p))$.
A cusp of label zero would force $B_p=1$, and hence
$1-t_p=1$ by full local Kummer injectivity.
A cusp of label one would force $t_p=1$.
At infinity the regular coordinate $t/(t-1)$ has
class one, forcing $-t_p/(1-t_p)=1$.
All three possibilities contradict $t_p\notin\{0,1\}$.
The local section is therefore the ordinary section
at $t_p$.

Globally $B=4A_{1/2}^2$ is a square, so (A.107.2) lifts
the section to $\pi_1(V,R)$.
Its finite local lift fixes a sheet of the degree-two
fibre above $t_p$, and hence is the section of an
actual point above $t_p$.  The two points are
$T_p$ and $1-T_p$.
Let $C$ denote the complete global class of
$T-1/2$ for this lifted section.
Equation (A.107.3) gives $C^2=A_{1/2}^2$.
The square kernel in $\widehat{\Q^\times}$ is
precisely $\{\pm1\}$: all its valuation coordinates
vanish because multiplication by two is injective
on $\widehat\Z$, and the remaining global valuation
kernel is $\{\pm1\}$, as proved in A.97.
Thus
\[
                              C=\pm A_{1/2}
                  \quad\hbox{with one global sign}.
 \tag{A.107.7}
\]
Apply the single deck involution if necessary to
make the sign positive.
At every finite prime, full local Kummer injectivity
now forces the lift's value of $T-1/2$ to equal
$T_p-1/2$.  Its actual centre is consequently $T_p$,
with the same deck choice at every finite prime.
The other two old classes have the local values
in (A.107.5); the global localization injection of A.76
makes them exactly $A_0,A_1$.

At the real place an ordinary point section lifts
to an ordinary real point: its fixed sheet is a
real point of the finite etale fibre.
A lifted cuspidal section is again cuspidal.
Indeed, the inertia orbits on the cover's geometric
sheets index the boundary points of the proper
normalization.  The lifted section fixes a sheet
and therefore its orbit; this gives a real boundary
point, and its decomposition group is the appropriate
lift of the cusp packet.
Thus the lifted global section satisfies the same
Selmer condition at the real place.
This proves surjectivity in (A.107.6).

For injectivity, two sections of the subgroup
conjugate in $\Pi$ are either conjugate in $\Lambda$
or become so after the deck involution: this is
the two-coset decomposition $\Pi/\Lambda=C_2$.
The second alternative sends $A_{1/2}$ to
$-A_{1/2}$ by (A.107.4).
These are distinct global classes, since $-1$ is
nontrivial in $\widehat{\Q^\times}$.
It cannot identify two classes on the left side
of (A.107.6).  This proves injectivity.
The pushforward itself preserves ordinary and
cuspidal local sections by the actual point and
boundary maps of (A.107.1), completing the proof of
the displayed bijection.

\paragraph{The finite geometric quotient maps.}
The old group $\Lambda$ is free profinite of rank three.
Define its characteristic subgroups
\[
 H_4=\ker\{\Lambda\longrightarrow\Lambda^{\rm ab}/4\},
                 \qquad K_4=\overline{H_4^2[H_4,H_4]}.
 \tag{A.107.8}
\]
The first subgroup is characteristic in $\Lambda$,
and the second is characteristic in $H_4$.
Hence $K_4$ is characteristic in $\Lambda$.
Because $\Lambda$ is normal in $\Pi$, conjugation
by every element of $\Pi$ preserves $K_4$.
The arithmetic action preserves it as well.
There are therefore actual arithmetic quotient maps
and an exact finite geometric sequence
\[
 \begin{gathered}
 \pi_1(V,R)/K_4\lhook\joinrel\longrightarrow
                      \pi_1(U,t_0)/K_4,\\
 1\longrightarrow\Lambda/K_4
       \longrightarrow\Pi/K_4
       \longrightarrow C_2\longrightarrow1.
 \end{gathered}
 \tag{A.107.9}
\]
The mod-four cover of $V$ has degree $4^3=64$.
Nielsen--Schreier gives
$\operatorname{rank}H_4=1+64(3-1)=129$.
Consequently
\[
 |H_4/K_4|=2^{129},\qquad
 |\Lambda/K_4|=2^{135},\qquad
                         |\Pi/K_4|=2^{136}.
 \tag{A.107.10}
\]
This uses the same old finite quotient as A.93;
normality under the larger group was proved above,
so no replacement by an unspecified core is needed.

The commutator inclusions also supply explicit
solvable quotient maps.  Since $\Pi'\subset\Lambda$,
one has
\[
 \Pi''\subset\Lambda'\subset H_4,\qquad
                  \Pi'''\subset[H_4,H_4]\subset K_4.
 \tag{A.107.11}
\]
Thus the actual third-derived quotient of the
tripod group maps onto $\Pi/K_4$.
The quotient $\Pi/K_4$ is a finite two-group.
Its abelianization has order at least four:
$K_4\subset\overline{\Pi^2[\Pi,\Pi]}$ and the
free rank-two group $\Pi$ has maximal elementary
abelian two-quotient $C_2^2$.
For a finite two-group of nilpotency class $c$,
its lower central series consequently gives order
at least $2^{c+1}$.
Equation (A.107.10) therefore gives the concrete bound
$c\le135$, and the actual quotient map
\[
                         \Pi/\Gamma_{136}\Pi
                                     \twoheadrightarrow\Pi/K_4.
 \tag{A.107.12}
\]
The point-induced arithmetic actions preserve this
map, so it also gives the corresponding arithmetic
quotient.
\end{proof}

The comparison (A.107.6) and the finite maps
(A.107.9)--(A.107.12) are proved independently of an
arithmetic obstruction example.
A.134 supplies an explicit completed harmonic triple
with a nonzero Brauer obstruction, and A.137 proves
its actual finite-quotient obstruction on the tripod.
The latter statement also computes this quotient's
exact class seventeen, sharpening (A.107.12) to the
map from $\Pi/\Gamma_{18}\Pi$.
